\documentclass{ltjsarticle}
\title{2次方程式の解の公式の導出}
\author{数学 太郎}
\usepackage{amsmath}
\begin{document}
\maketitle
2次方程式$ax^2 + bx + c = 0\,(a\neq 0)$の解は以下のように係数の四則演算と根号を使って表せます。これを2次方程式の\textbf{解の公式}と呼びます。
\[ x = \frac{-b \pm \sqrt{b^2 -4ac}}{2a} \]

\section{導出}
2次方程式の解の公式は、以下のように平方完成を使うと導出できます。
\begin{align*}
ax^2 + bx + c &= a\left(x^2 + \frac{b}{a}x\right) + c\\
&= a\left(x^2 + \frac{b}{a}x + \frac{b^2}{4a^2} - \frac{b^2}{4a^2}\right) + c\\
&= a\left( x + \frac{b}{2a} \right)^2 - \frac{b^2}{4a} + c\\
&= a\left( x + \frac{b}{2a} \right)^2 - \frac{b^2-4ac}{4a}
\end{align*}
よって、
\begin{align*}
a\left( x + \frac{b}{2a} \right)^2 - \frac{b^2-4ac}{4a} &= 0\\
a\left( x + \frac{b}{2a} \right)^2  &= \frac{b^2-4ac}{4a}\\
\left( x + \frac{b}{2a} \right)^2  &= \frac{b^2-4ac}{4a^2}
\end{align*}
両辺の根号をとって変形すると、
\begin{align*}
x + \frac{b}{2a} &= \pm \frac{\sqrt{b^2-4ac}}{2a}\\
x &= -\frac{b}{2a}\pm \frac{\sqrt{b^2-4ac}}{2a}\\
x &= \frac{-b \pm \sqrt{b^2-4ac}}{2a}\\
\end{align*}
\section{平方完成の積分への応用}
平方完成を応用すると、次の形の積分を計算することができます。
\[ \int \frac{dx}{ax^2 + bx + c} \]

例として、次の積分を計算してみましょう。
\[ \int \frac{dx}{x^2 + x2 + 2} \]
分母を平方完成します。
\begin{align*}
\int \frac{dx}{x^2 + 2x + 2} &= \int \frac{dx}{(x^2 + 2x + 1 - 1) + 2}\\
&= \int \frac{dx}{(x^2 + 2x + 1) - 1 + 2}\\
&= \int \frac{dx}{(x+1)^2 + 1}\\
&= \arctan{(x+1)} + C
\end{align*}
ただし、$C$は積分定数です。
\end{document}

